\section{Giới thiệu}
Khai thác tập phổ biến (frequent itemset mining, FIM) tìm các tập item cùng xuất hiện
trong ít nhất $\sigma_{\min}$ giao dịch của một cơ sở dữ liệu, trong đó $\sigma_{\min}$ là
ngưỡng hỗ trợ (minimum support, viết tắt minsup khi tính theo tỉ lệ). Đây là bước đầu của khai
thác luật kết hợp và là nền của nhiều bài toán khai thác mẫu. Số tập phổ biến có thể tăng
rất nhanh khi ngưỡng giảm, nên cách biểu diễn dữ liệu và cách đếm support quyết định thuật
toán chạy được ở ngưỡng nào.

Eclat~\cite{zaki2000} biểu diễn cơ sở dữ liệu theo chiều dọc: mỗi item đi kèm tidset,
tức tập mã giao dịch chứa nó. Support của một tập item bằng kích thước giao của các
tidset, và Eclat tìm kiếm theo chiều sâu trên các lớp tương đương (equivalence class)
chia theo tiền tố. Nhóm cài đặt lại thuật toán này từ bài báo gốc bằng Julia, chỉ dùng thư
viện chuẩn, so với Eclat và PrePost+ của SPMF~\cite{spmf2016}, và dùng nó cho hai ứng dụng
thực tế. Bài báo gồm ba phần chính theo đề:
\begin{itemize}
\item Mục~\ref{sec:implementation} trình bày bản cơ bản bám sát bài báo gốc và bản tối ưu
(bitset, thứ tự duyệt theo support tăng dần, tỉa phần mở rộng hoàn hảo, đếm cặp và danh
sách tid cho dữ liệu thưa), cùng chứng minh tính đúng, cách kiểm thử và mức cải thiện đo
được của từng kỹ thuật tối ưu.
\item Mục~\ref{sec:evaluation} nêu giao thức đo chung (Mục~\ref{sec:protocol}), đối chiếu
toàn bộ tập mục với SPMF (Mục~\ref{sec:exact}), thời gian theo minsup
(Mục~\ref{sec:minsup}), số tập phổ biến theo minsup (Mục~\ref{sec:output-size}), bộ nhớ
đỉnh (Mục~\ref{sec:memory}), khả năng mở rộng theo kích thước dữ liệu
(Mục~\ref{sec:scalability}), ảnh hưởng của độ dài giao dịch (Mục~\ref{sec:width}), và
hai hướng tối ưu tiếp theo (Mục~\ref{sec:future}).
\item Mục~\ref{sec:applications} áp dụng Eclat vào phân tích giỏ hàng trên Online Retail
II (Mục~\ref{sec:basket}) và khai thác mẫu hình ảnh trên Semeion
(Mục~\ref{sec:visual-patterns}).
\end{itemize}

Hai nguyên tắc chi phối cách trình bày. Thứ nhất, mọi so sánh hiệu năng chỉ dùng các lần
chạy có đầu ra khớp hoàn toàn, về tập mục và support, với một lần chạy Eclat của SPMF
độc lập, vì hai chương trình có thể cho cùng số tập mục mà gán support khác nhau. Thứ hai,
các số đo chính gắn với máy và revision mã nguồn đã tạo ra chúng, và mọi bước ngẫu nhiên
dùng seed 42. Repository của dự án chứa mã nguồn, cấu hình, mã băm dữ liệu và lệnh tái lập
của mọi thực nghiệm.
